% ============ 第2章 §2.4 Poisson-Stieltjes积分表示(原书 pp.132-134) ============
\section{Poisson--Stieltjes 积分表示}\label{sec:2.4}

本节证明属于 $h^{1}$ 的函数可以用 Poisson--Stieltjes 积分表示.因此证明属于 $h^{1}$ 的函数都拥有径向极限.

\begin{theorem}\label{thm:2.4.1}
若 $u \in h^{1}$,则存在唯一的实测度 $\mu \in M[-\pi, \pi]$ 满足 $\mu(\{-\pi\}) = \mu(\{\pi\})$,且
\begin{equation*}
u(re^{i\theta}) = \frac{1}{2\pi}\int_{-\pi}^{\pi}P(r, \theta - t)\dd\mu(t).
\end{equation*}
\end{theorem}

\begin{proof}
令 $r_{n} = 1 - \frac{1}{n}$ 且 $u_{n}(z) = u(r_{n}z)$,则 $u_{n}(z)$ 在 $|z| < 1/(1 - \frac{1}{n})$ 上调和,故如 2.2 节所示,
\begin{equation*}
u_{n}(z) = \frac{1}{2\pi}\int_{-\pi}^{\pi}P(r, \theta - t)u_{n}(e^{it})\dd t,
\end{equation*}
可以用 Poisson 积分表示.由于 $u \in h^{1}$,
\begin{align*}
\|u_{n}\|_{1} & = \int_{-\pi}^{\pi}|u_{n}(e^{it})|\dd t \leq \sup_{0 \leq r < 1}\int_{-\pi}^{\pi}|u(re^{it})|\dd t \\
& = 2\pi\sup_{0 \leq r < 1}M_{1}(r, u) < \infty,
\end{align*}
所以 $\{u_{n}\}$ 属于 $L^{1} = L^{1}[-\pi, \pi]$ 的半径为 $C$ 的球.这里 $C = 2\pi\sup_{0 \leq r < 1}M_{1}(r, u)$.

如 2.2 节所注意到的,$L^{1}$ 可以等距嵌入 $M = M[-\pi, \pi]$.$M$ 是 $[-\pi, \pi]$ 上连续函数的全体 $C[-\pi, \pi]$ 的对偶空间.因 $M$ 的单位球在弱$^{*}$拓扑下紧,故存在 $\{u_{n}\}$ 的子列 $\{u_{n_{j}}\}$ 与 $\mu \in M$,使得 $u_{n_{j}}$ 弱$^{*}$收敛于 $\mu$.此时 $\mu(\{-\pi\}) = \mu(\{\pi\})$.因此,固定 $r$ 与 $\theta$,因 $P(r, \theta - t) \in C[-\pi, \pi]$,
\begin{equation*}
\lim_{j \rightarrow \infty}\int_{-\pi}^{\pi}P(r, \theta - t)u_{n_{j}}(e^{it})\dd t = \int_{-\pi}^{\pi}P(r, \theta - t)\dd\mu(t)
\end{equation*}
成立.对每个 $j$,
\begin{equation*}
u_{n_{j}}(re^{i\theta}) = \frac{1}{2\pi}\int_{-\pi}^{\pi}P(r, \theta - t)u_{n_{j}}(e^{it})\dd t
\end{equation*}
且 $\lim_{j \rightarrow \infty}u_{n_{j}}(re^{i\theta}) = u(re^{i\theta})$,所以
\begin{equation*}
u(re^{i\theta}) = \frac{1}{2\pi}\int_{-\pi}^{\pi}P(r, \theta - t)\dd\mu(t)
\end{equation*}
成立.$\mu$ 的唯一性是命题~\ref{prop:2.2.3} 的结果.
\end{proof}

定理~\ref{thm:2.4.1} 中当 $u \in h^{p}$($1 < p \leq \infty$)时,可以取 $d\mu = f(t)\dd t$ 且 $f \in L^{p}$.这是因为,设 $1/p + 1/q = 1$,则 $L^{p} = (L^{q})^{*}$,可以像定理~\ref{thm:2.4.1} 的证明那样证明.

若 $f \in H^{1}$ 且 $f = u + iv$,则 $u, v \in h^{1}$,所以 $f$ 可以用复测度作 Poisson--Stieltjes 积分.2.12 节将证明该测度关于 Lebesgue 测度绝对连续.

\begin{corollary}\label{cor:2.4.2}
若 $u \in h(D)$ 在 $D$ 上为正,则存在正测度 $\mu \in M[-\pi, \pi]$ 满足 $\mu(\{-\pi\}) = \mu(\{\pi\})$,且
\begin{equation*}
u(re^{i\theta}) = \frac{1}{2\pi}\int_{-\pi}^{\pi}P(r, \theta - t)\dd\mu(t).
\end{equation*}
\end{corollary}

\begin{proof}
$u \in h^{1}$.因为
\begin{equation*}
\int_{-\pi}^{\pi}|u(re^{i\theta})|\dd\theta = \int_{-\pi}^{\pi}u(re^{i\theta})\dd\theta = 2\pi u(0).
\end{equation*}
由定理~\ref{thm:2.4.1},$u$ 可以由 $\mu$ 作积分表示,但 $\mu$ 是由 $u$ 定义的 $u_{n_{j}}$ 的弱$^{*}$极限.由 $u_{n_{j}}(z) > 0$($z \in \overline{D}$)知 $\mu$ 是正测度.
\end{proof}
